\section{Feasible Shifts and Volume Distortion}
\label{sec:feasible-shifts}

The geometric arguments will need to raise or lower heights slowly across
a corridor while retaining an event of uniformly positive probability.
Simply adding the prescribed shift can violate a nearly saturated edge
constraint. We correct the shift by averaging it over clusters of such
edges, so that the two corrected endpoint shifts agree at every exact
edge contact. A
determinant identity for vector fields tangent to the height polytope
then controls the volume lost by this height-dependent correction.
The resulting proposition transfers any event of probability bounded
below to a feasible shifted event of probability bounded below; the
geometric application must supply that initial event.

Constraint-preserving perturbations and their volume cost also underlie
Richthammer's deformed translations for hard-core particles
\cite[Sections~5.4 and~6.5--6.7]{Richthammer} and the addition algorithm
of Mi\l o\'s and Peled for hard-constraint surfaces
\cite[Sections~1.1.2 and~2]{MilosPeled}. The threshold-averaged component
construction below differs from those algorithms and supplies its own
tangential-vector-field determinant identity.

All heights in this section are in lattice units, and the law is
normalized Lebesgue measure on a positive-volume conditional height
polytope. Conditioning specifies signs along oriented strands and may
include the bounded fixed data allowed in
Corollary~\ref{f:conditioned-sparse}. Sites with specified signs will
be left unchanged by the map. The sparse edge estimates of that
corollary apply also to zero-pinned laws by taking no sign strings.

\subsection{The cluster consequence of sparseness}

Fix a bounded-degree lattice graph and a set of at most $C n^2$ vertices.
For a sufficiently small fixed $s_*>0$, join neighboring vertices when
\[
 |h_x-h_y|\geq 1-s_*.
\]
Ignore edges whose two endpoints are fixed. The sparse edge estimate of
Corollary~\ref{f:conditioned-sparse}, including its zero-pin specialization,
implies that the size $H_x$ of the resulting component containing $x$
satisfies
\begin{equation}\label{eq:shift-cluster-moments}
 \sup_x \mathbb E H_x^k\leq C_k,\qquad k<\infty.
\end{equation}
The constants are uniform over the conditioning systems covered by
that corollary. To see the implication, use its overlapping fixed-box
cover: the central-edge tests include every edge with a nonfixed endpoint,
also near fixed pins. A nontrivial component can occupy only bad boxes,
apart from a bounded enlargement at its endpoints, since a good box has
slack on every edge incident to its central box. A connected collection
of $m$ occupied boxes contains a separated subcollection of at least
$m/C$ boxes. Its probability is at most $\rho^{m/C}$. Taking $\rho$
smaller than the reciprocal of the lattice-animal counting constant
gives an exponential tail for the number of occupied boxes, hence
\eqref{eq:shift-cluster-moments}.
Deleting edges or restricting the graph preserves this conclusion.
In particular, for any fixed $\alpha>0$,
\begin{equation}\label{eq:shift-largest-cluster}
 \mathbb P\left[\max_x H_x>n^\alpha\right]
 \leq C_k n^{2-\alpha k}=o(1)
\end{equation}
after choosing $k>2/\alpha$. No convergence rate from homogenization enters
this assertion.

\subsection{A determinant identity on polytopes}

\begin{lemma}[Tangential vector fields]\label{lem:shift-divergence}
Let $P\subset\mathbb R^N$ be a compact full-dimensional polytope, with
normalized Lebesgue measure. Suppose that $u:P\to\mathbb R^N$ is Lipschitz
and tangent to each facet of $P$. Then
\begin{equation}\label{eq:shift-divergence}
 \mathbb E\operatorname{div}u=0,\qquad
 \mathbb E(\operatorname{div}u)^2
 =\mathbb E\operatorname{tr}\bigl((Du)^2\bigr)
 \leq\mathbb E\|Du\|_{\mathrm{HS}}^2.
\end{equation}
\end{lemma}

\begin{proof}
Write $\ell_i\geq0$ for the normalized affine slack of facet $F_i$.
There is a finite constant $C_i$, depending here on the polytope, such that
$\operatorname{dist}(x,F_i)\leq C_i\ell_i(x)$ for $x\in P$. Indeed, express
$x$ as a convex combination of vertices, fix $v_0\in F_i$, and replace
each vertex $v\notin F_i$ in that combination by $v_0$; use the maximum of
$|v-v_0|/\ell_i(v)$ over these finitely many vertices.
Tangency and Lipschitz continuity give
\[
 |\nabla\ell_i\cdot u(x)|
 \leq C_i\operatorname{Lip}(u)\ell_i(x).
\]
Consequently, for sufficiently small positive or negative $t$,
$T_t=I+tu$ maps $P$ into itself and maps $\partial P$ into $\partial P$.
It is injective when $|t|\operatorname{Lip}(u)<1$. To prove surjectivity,
fix $y\in P$. The map $x\mapsto\Pi_P(y-tu(x))$, where $\Pi_P$ is
Euclidean projection, is a contraction of $P$ and has a fixed point $x$.
Put $n=y-tu(x)-x$. The projection inequality says that $n$ is an
outward normal to $P$ at $x$, so $n\cdot(y-x)\le0$. Tangency to every
facet containing $x$ gives $n\cdot u(x)=0$, because the normal cone
is generated by those facet normals. Consequently
$0\ge n\cdot(y-x)=|n|^2$, so $n=0$ and $T_t(x)=y$.

The Lipschitz area formula now gives
\[
 \int_P \det(I+tDu(x))\,dx=|P|.
\]
For the piecewise affine field constructed below, this identity is
elementary: subdivide $P$ into its finitely many affine regions,
apply the linear determinant formula on each region, and sum.
Injectivity makes the image interiors disjoint, while their boundaries
have zero volume; surjectivity makes their union all of $P$.
The determinant is positive for these $t$. Comparing the linear and
quadratic coefficients of this polynomial identity proves
\eqref{eq:shift-divergence}, since the quadratic coefficient is
$((\operatorname{div}u)^2-\operatorname{tr}((Du)^2))/2$.
Finally $|\operatorname{tr}(A^2)|\leq\|A\|_{\mathrm{HS}}^2$.
Thus no integration by parts across corners or nonsmooth interfaces is
required. The polytope-dependent restriction on $t$ has disappeared from
the resulting identity.
\end{proof}

\subsection{Construction and quantitative bounds}

Let $V_n$ contain the sites on which the prescribed shift may be nonzero.
Assume $|V_n|\leq Cn^2$, and adjoin all their neighbors. Declare all adjoined
vertices, all fixed sites, and all sites carrying sign or additional wall
conditions to be grounded. Include graph edges with at least one
nongrounded vertex, and omit edges with two grounded endpoints.
Grounding means only that the shift is zero there; a grounded height need
not be fixed under the original law.

Let $f=f_n$ be a deterministic prescription, zero at every grounded site
and outside $V_n$, such that
\begin{equation}\label{eq:shift-prescription}
 |f_x|\leq M,\qquad |f_x-f_y|\leq L/n
 \quad\hbox{for every lattice edge }xy.
\end{equation}
The constants $L,M$ may be arbitrarily large fixed constants. They do not
depend on $n$. Choose fixed
$0<s_0<s_1<s_2<s_*$, and put $\Delta=s_1-s_0$.
For an edge $e=xy$, put $t_e(h)=1-|h_x-h_y|$.
For $s\in[s_0,s_1]$, let $C_s(x)$ be the component of $x$ using edges
with $t_e(h)\leq s$, and define
\[
 v_s(h)_x=
 \begin{cases}
  |C_s(x)|^{-1}\displaystyle\sum_{y\in C_s(x)}f_y,
       & C_s(x)\hbox{ has no grounded vertex},\\
  0,   & C_s(x)\hbox{ has a grounded vertex}.
 \end{cases}
 \qquad
 u(h)=\Delta^{-1}\int_{s_0}^{s_1}v_s(h)\,ds.
\]
Set $u=0$ outside the local graph. Let $H_x$ now denote the size of the
$s_2$-component containing $x$, and write $H_{\max}=\max_x H_x$.
The definitions give the deterministic approximation estimate
\begin{equation}\label{eq:shift-error}
 |u(h)_x-f_x|\leq L H_x/n.
\end{equation}
For an ungrounded component this is the oscillation bound for its average.
For a grounded component, join $x$ to a grounded vertex along its component
and use $f=0$ at the latter vertex.

The vector field $u$ is globally Lipschitz and piecewise affine on the
finite-dimensional height polytope. Indeed, subdivide into the finitely many
regions on which the signs of edge differences and the ordering of
$s_0,s_1,t_e(h)$ are fixed. On each such region the component partitions
between consecutive thresholds are fixed and the lengths of the
corresponding integration intervals are affine. The resulting formulas
agree continuously at common boundaries. A direct continuity bound follows
also by noting that changed edge statuses can occur only for $s$ between
$t_e(h)$ and $t_e(\widetilde h)$ and summing these interval lengths.

At an edge facet $|h_x-h_y|=1$, its endpoints belong to the same component
for every $s\in[s_0,s_1]$, so $u_x=u_y$. At every wall or fixed site,
$u_x=0$. Thus $u$ is tangent to every facet of the height polytope, as
required in Lemma~\ref{lem:shift-divergence}.

Here is an explicit derivative bound. When $s$ passes $t_e$, write
$J_e=v_{t_e+}-v_{t_e-}$ for the change of the component prescription.
If the edge is redundant, $J_e=0$. Otherwise it merges two components;
$J_e$ is supported on their union, and
$|(J_e)_x|\leq L H_x/n$ at each affected vertex. This remains true when
one component is grounded. Almost everywhere,
\begin{equation}\label{eq:shift-derivative}
 Du=-\Delta^{-1}
       \sum_{\{e:s_0<t_e<s_1\}} J_e\otimes\nabla t_e.
\end{equation}
Equal thresholds can be ordered by a fixed tie rule; simultaneous mergers
then telescope to the same formula. Each $\nabla t_e$ is supported on
the two endpoints of $e$ and has entries of magnitude at most one.
Hence each column of \eqref{eq:shift-derivative} receives contributions
from at most the graph degree many edges. Matrix entries connecting
different $s_1$-components vanish. On a component of size $H$,
each entry has magnitude at most $C L H/n$. Summing entries, and then
squares, yields
\begin{equation}\label{eq:shift-matrix-bounds}
 \|Du\|_{\mathrm{op}}\leq C L H_{\max}^2/n,\qquad
 \|Du\|_{\mathrm{HS}}^2\leq C L^2n^{-2}\sum_x H_x^3.
\end{equation}
All constants may depend on the fixed thresholds. In particular,
\eqref{eq:shift-cluster-moments} and
Lemma~\ref{lem:shift-divergence} imply
\begin{equation}\label{eq:shift-mean-bounds}
 \mathbb E\|Du\|_{\mathrm{HS}}^2\leq C_L,\qquad
 \mathbb E(\operatorname{div}u)^2\leq C_L.
\end{equation}
These expectations are under the original conditional law; no derivative
bound is assumed after an additional exploration.

\subsection{Feasibility, injectivity, and volume}

\begin{proposition}[Feasible shifts]\label{g:shifts}
Under the sparse-box estimates \ref{f:sparse} and
\ref{f:conditioned-sparse} in the applicable conditioning class, let $f_n$ satisfy
\eqref{eq:shift-prescription}. For each $p_0>0$ there are $c>0$ and $n_0$,
uniform over the allowed conditioning systems, with the following property.
For $n\geq n_0$ and every measurable event $E$ with $\mathbb P(E)\geq p_0$,
there is a measurable $E'\subset E$ with $\mathbb P(E')\geq p_0/2$ such that
$T(h)=h+u(h)$ is injective on $E'$, has feasible image, and
\[
 \sup_{h\in E'}\|T(h)-h-f_n\|_\infty=o(1),\qquad
 \mathbb P(T(E'))\geq c.
\]
The statement also holds for a fixed finite product of independent
conditional surfaces.
\end{proposition}

\begin{proof}
Take $\alpha=1/8$ and
$G_n=\{H_{\max}\leq n^\alpha\}$, using the enlarged threshold $s_2$.
By \eqref{eq:shift-largest-cluster},
$\mathbb P(G_n^c)=o(1)$ uniformly.
If $h\in G_n$, then \eqref{eq:shift-error} gives
$\|u-f\|_\infty\leq Ln^{\alpha-1}$.
For an edge with $t_e\leq s_0$, its two shifts agree exactly.
For every other edge the original slack is greater than $s_0$, whereas
\[
 |u_x-u_y|\leq L(2n^\alpha+1)/n<s_0
\]
for large $n$. Thus $h+u(h)$ is feasible. Walls are preserved because
their sites are grounded.

To prove injectivity, suppose $h,k\in G_n$ and $h+u(h)=k+u(k)$.
The prescribed $f$ is the same for both fields, so
\begin{equation}\label{eq:shift-collision}
 \|h-k\|_\infty\leq 2Ln^{\alpha-1}.
\end{equation}
The whole segment $h_t=(1-t)h+tk$ belongs to the convex polytope.
If an edge has $t_e(h_t)\leq s_1$, then
\[
 t_e(h)\leq s_1+2\|h_t-h\|_\infty
 \leq s_1+4Ln^{\alpha-1}<s_2
\]
for large $n$. Every $s_1$-component anywhere along this segment is
therefore contained in an $s_2$-component of $h$ and has size at most
$n^\alpha$. Apply the component proof of
\eqref{eq:shift-matrix-bounds} along the segment to get
\[
 \|u(k)-u(h)\|_2
 \leq C L n^{2\alpha-1}\|k-h\|_2.
\]
For a segment lying on a subdivision boundary, the same conclusion follows
by restricting the continuous piecewise-affine formula to that segment;
the tangential derivatives agree with limits from adjacent pieces.
Since the coefficient tends to zero, the collision equation forces $h=k$.
Notice that this argument does not presume that $G_n$ is convex.

On $G_n$, $\|Du\|_{\mathrm{op}}=o(1)$, so for large $n$ the determinant
of $I+Du$ is positive and
\begin{equation}\label{eq:shift-logdet}
 \log\det(I+Du)\geq
       \operatorname{div}u-\|Du\|_{\mathrm{HS}}^2.
\end{equation}
Indeed, use the power series for $\log(I+A)$ and
$|\operatorname{tr}(A^j)|
\leq\|A\|_{\mathrm{HS}}^2\|A\|_{\mathrm{op}}^{j-2}$ for $j\geq2$;
the remainder is at most $\|A\|_{\mathrm{HS}}^2$ when
$\|A\|_{\mathrm{op}}\leq1/2$.
By \eqref{eq:shift-mean-bounds}, choose a fixed $K=K(p_0,L)$ sufficiently
large that
\[
 \mathbb P\bigl[
 |\operatorname{div}u|>K
 \ \hbox{or}\ \|Du\|_{\mathrm{HS}}^2>K\bigr]<p_0/4.
\]
For sufficiently large $n$, $\mathbb P(G_n^c)<p_0/4$ also. Intersecting
$E$ with the complementary good events gives $E'$ of probability at least
$p_0/2$, and \eqref{eq:shift-logdet} gives
$\det DT\geq e^{-2K}$ almost everywhere on $E'$.
The area formula and injectivity now imply
\[
 \mathbb P(T(E'))
 =\frac1{|P|}\int_{E'}\det DT(h)\,dh
 \geq e^{-2K}p_0/2.
\]
This proves the claim. For finitely many independent replicas use the
product polytope and the block-diagonal map, or intersect the corresponding
good events and multiply the determinant lower bounds.
\end{proof}

The construction also preserves the direction of a reinforcing shift.
If $f$ is nonnegative on an $s_1$-component, its corrected shift is
nonnegative there, including when it is grounded. If positive and negative
prescription regions have separation of order $n$, the components on $G_n$
cannot connect them. This observation permits opposite prescribed signs
in separated corridors. A crossing application must still establish its
own initial event $E$ and verify that an $o(1)$ error in the prescribed
shift leaves its required strict inequalities intact.
